\chapter{合作、竞争与多层次选择}

\begin{kaichang}
厨房里那罐蜂蜜，是昆虫世界最惊人的工程之一。一汤匙蜂蜜，大约需要十几只工蜂忙活一辈子：一只工蜂一生只能酿出约十二分之一茶匙的蜜。它们飞出巢门，访问上百万朵花，把花蜜运回巢中，扇动翅膀扇干水分，封进蜡格——而这一切劳动的果实，大部分它们自己根本吃不到。

更奇怪的事还在后面。工蜂全部是雌性，却几乎从不产卵；遇到危险，它们会毫不犹豫地亮出螫针——蜜蜂的螫针带着倒钩，扎进皮肤后会连同一部分内脏一起被扯掉，蜇人的工蜂几小时内就会死去。\textbf{放弃生育，为集体送命}。按照第 76 章的逻辑，自然选择应该偏爱那些拼命留下自己后代的个体。一个“只顾别人”的性状，凭什么不被淘汰，反而流传了几千万年？

先别急着翻答案。拿一张纸，写下你的猜想。这一章结束时，我们会回到这罐蜂蜜——并且你会看到，谜底不仅解释了蜜蜂，还解释了你身体里的每一个细胞，甚至解释了癌症。
\end{kaichang}

\section{合作遍地都是，竞争也遍地都是}

先把镜头拉远，看看自然界里合作到底普遍到什么程度：

\begin{itemize}[itemsep=2pt]
\item 蚂蚁在植物茎上“放牧”蚜虫：蚜虫吸食植物汁液，排出含糖的蜜露喂蚂蚁；蚂蚁则驱赶瓢虫，保护蚜虫。一方出糖，一方出兵。
\item 岩石上灰绿色的地衣，不是单一生物，而是真菌和藻类（或蓝细菌）的合伙公司：藻类光合作用制造糖，真菌提供住所、水分和矿物质。拆散它们，双方在荒野里都很难独活。
\item 大豆、豌豆、花生的根上长着一粒粒小瘤，里面住着根瘤菌。细菌把空气中的氮气变成植物能用的含氮养料，植物则把光合作用生产的糖分给细菌。
\item 你自己的肠道里住着以万亿计的细菌，它们帮你分解膳食纤维、合成某些维生素；你则提供温暖、湿润、永不缺席的三餐。
\end{itemize}

而同一个舞台上，竞争同样无处不在：蜂箱里，上万只蜜蜂争夺同一片花丛；田地里，杂草和庄稼争夺阳光、水分和氮；你的肠道里，不同细菌菌株争夺肠壁上每一处落脚点。合作与竞争不是两个对立的世界，它们是同一套规则的两种结局。这套规则你已经认识了，用第 76 章的正式说法：\textbf{可遗传的变异，加上差异繁殖成功，等于演化}。

但前几章我们一直把“个体”当作被选择的单位：长脖子的个体活得好，深色的个体藏得好。蜜蜂的悬案恰恰戳在这个默认假设上：如果选择的单位是个体，工蜂的利他行为完全讲不通。所以本章真正的问题是——\textbf{自然选择到底在哪一层发生}？是基因？是个体？还是群体？

\section{亲缘选择：帮亲戚，就是帮自己的副本}

1964 年，英国生物学家汉密尔顿给出了一个干净得让人不安的回答。他的想法用一句话就能说清：\textbf{基因不在乎它住在谁的身体里}。

回想第 69 章：有性生殖时，父母各出一半基因组成受精卵。所以你和你兄弟姐妹平均共享 $1/2$ 的基因；你和你的表亲平均共享 $1/8$；你和你父母、子女之间也是 $1/2$。这个共享比例叫做\emph{亲缘系数}，记作 $r$。现在想象这样一个场景：一只动物面前摆着两个选择——自己逃走（保住自己这套基因），或者发出警报引来天敌、大概率送命（但救下它身边的一群亲戚）。警报声让自己的这套基因消失了，可亲戚们身体里还躺着大量相同的基因副本，安然无恙。从“基因的平均命运”来算账，警报完全可能是笔好买卖。

\begin{qiaoliang}
汉密尔顿把这笔账写成了一条不等式，后人叫它\emph{汉密尔顿法则}：
\[
rB > C
\]
其中 $C$ 是利他行为给自己繁殖带来的损失（cost），$B$ 是给受助者繁殖带来的好处（benefit），$r$ 是两者的亲缘系数。只要不等式成立，“帮助亲戚”的基因就倾向于在群体中增加。

算一笔具体的账：牺牲自己（$C=1$，损失全部未来繁殖），能救下多少个表亲才“够本”？表亲的 $r=1/8$，每个表亲获救的收益按 $B=1$ 计，则 $\frac{1}{8} \times N > 1$，需要 $N>8$。所以有句半开玩笑的名言（出自比汉密尔顿更早的遗传学家霍尔丹）：“我愿为两个亲兄弟，或者八个表兄弟，跳进河里。”——兄弟 $r=1/2$，两个就够；表亲 $r=1/8$，得凑八个。数字当然是极度简化的（现实中 $B$ 和 $C$ 很少恰好等于 1），但它抓住了要害：\textbf{帮助越亲的亲属、代价越小、收益越大，利他行为就越容易被演化保留}。这就是\emph{亲缘选择}。
\end{qiaoliang}

现在回头看蜜蜂，还有一层更妙的结构。蜂类、蚁类是\emph{单双倍体}：受精的卵发育成雌性（工蜂和蜂王），未受精的卵直接发育成雄蜂。雌蜂从父亲那里得到一整套基因（父亲只有一套可给），从母亲那里得到一半。算下来，一只工蜂与它的姐妹平均共享 $3/4$ 的基因——比它与亲生女儿能共享的 $1/2$ 还要高！也就是说，对一只工蜂而言，“帮蜂王妈妈多养一个妹妹”比“自己生一个女儿”在基因账上更划算。这曾是解释蜜蜂不育品级的经典答案。

不过要诚实地补一句：近几十年的研究发现，许多独居的蜂和蚁同样是单双倍体，而一些高度社会性的昆虫（比如白蚁）却不是。所以今天的科学家大多认为，$3/4$ 这个巧合不是关键，关键是汉密尔顿那条更一般的逻辑——\textbf{帮助亲属，等价于帮助自己基因的其他副本}。单双倍体只是把这个逻辑拧得更紧的一颗螺丝。科学观点会随证据修正，这个修正过程本身，就是第 76 章以来我们一直在演示的方法。

亲缘选择不只解释昆虫。地松鼠看到天敌时会发出尖叫报警，报警者更危险，但周围的亲属得以逃生；许多鸟类（如某些松鸦）里，年轻个体先不繁殖，留在父母的巢边帮忙喂养弟妹——“当帮手”同样是那笔亲缘账。

\section{互惠与共生：没有血缘，也能合伙}

如果合作只能靠血缘，那蚂蚁和蚜虫、真菌和藻类、你和肠道细菌就都没法解释了。非亲非故的合作，走的是另外两条路。

第一条路叫\emph{互惠}：这次我帮你，下次你帮我。它的经典例子是吸血蝙蝠。这种蝙蝠每夜必须吸到血，连续两三夜落空就可能饿死；而满载而归的个体，常常会把一部分血反刍给同一栖息群里没吃到东西的同伴——哪怕不是亲属。野外观察发现，反哺明显偏向“曾经喂过我”的个体。互惠要成立，需要三个条件，你可以当成 checklist：双方\textbf{反复相遇}（不是一锤子买卖）、能\textbf{认出对方}（防止只占便宜不回报的骗子白吃）、帮助一方的\textbf{代价小于}受助一方的收益（我血多喝不完，你快饿死，一口血对你价值巨大）。三条齐备，“记账式合作”就能稳住。

第二条路走得更远，叫\emph{共生}：两个物种的代谢直接咬合在一起，你缺的正是我多的。大豆根瘤是最好的例子，也是我们这一章的动手对象。

\begin{shiyan}
\textbf{观察豆科植物的根瘤}（安全的野外观察）

\textbf{材料}：一棵豆科植物（田边的三叶草、野豌豆、苜蓿，或花盆里的豆类幼苗）、小铲、一盆清水、放大镜、一次性手套、一张白纸。

\textbf{步骤}：1. 用小铲连根挖起整株植物（选幼小植株，尽量保全细根）；2. 在水中轻轻漂掉泥土，把根摊在白纸上；3. 用放大镜仔细找根上附着的小米粒大小的突起——那就是根瘤，和非豆科植物的根（比如顺便挖的一棵小草）对比；4. 用指甲轻轻掰开一个较大的根瘤，看内部的颜色。

\textbf{观察点}：根瘤有多少个、长在主根还是侧根上？掰开的根瘤内部往往呈淡红色——那是\emph{豆血红蛋白}，一种和动物血红蛋白相似的蛋白质，由植物和细菌合作制造，负责调节根瘤内的氧气浓度，保护固氮反应。

\textbf{安全提示}：戴手套操作，结束后用肥皂洗手；不认识的植物不要入口，观察完毕把植物放回原处或妥善处理。
\end{shiyan}

根瘤里发生的，是全书化学链上重要的一环。空气中的氮气 \chem{N_2} 极其稳定——两个氮原子之间是三键，拆开它要付出很高的能量代价（第 30 章前后讲化学键能量时算过这笔账）。根瘤菌体内的\emph{固氮酶}能完成这件事，但代价不菲：把一分子的 \chem{N_2} 还原成两分子氨（\chem{NH_3}），至少要消耗 16 个 ATP，而 ATP 的原料正是植物光合作用送来的糖。这笔交易里，植物出糖，细菌出固氮的化学功夫，双方都离不开对方。农业上种植豆科植物“养地”，利用的就是这条共生通道：收割后留在土壤里的含氮物质，能减少化肥的使用量。

\section{最成功的合作：把伙伴变成自己的一部分}

共生还能走得更极端：一方干脆搬进另一方体内，再也不分开。这不是科幻，而是你身体里已经发生过的历史——你每个细胞里的线粒体（第 58 章的“发电厂”），以及每片叶子里的叶绿体（第 54 章的光合车间），祖先都是自由生活的细菌。

\begin{zhengju}
1967 年，美国生物学家林恩·马古利斯（当时还姓萨根）把一篇论文投给学术期刊，核心主张惊世骇俗：真核细胞里的线粒体和叶绿体，起源于被吞进细胞、却没被消化、反而定居下来的细菌。这篇论文被\textbf{退了大约十五次稿}才发表，此后多年被主流学界当成异想天开。

但马古利斯不是在空想，她手里有一串可以被检验的观察。线粒体和叶绿体都有一层额外的膜（双层膜——内层像细菌自己的膜，外层像宿主吞食时包上去的“口袋”）；它们各自带着一小份环状 DNA，形态像细菌的染色体，而不像细胞核里的线状染色体；它们内部的核糖体尺寸更像细菌的，某些能抑制细菌核糖体的抗生素，也能抑制线粒体的核糖体；它们在细胞里以类似细菌二分裂的方式增殖，细胞造不出“从零开始”的线粒体——线粒体只能来自已有的线粒体。这些现象用“细胞器是吞进来的细菌”能统一解释，用别的假说解释起来处处别扭。

真正一锤定音的证据，要到 DNA 测序技术成熟之后才到来：把线粒体 DNA 的序列与各种细菌逐一比对，发现它与一类叫 $\alpha$-变形菌的细菌最接近；叶绿体的序列则指向蓝细菌。祖先的“户口”找到了。今天，内共生起源已是教科书共识；至于一些更细的问题——比如宿主与细菌的基因是怎样一步步“搬家”到细胞核里的、其他细胞结构是否也有内共生来源——仍在研究和争论之中，科学还没有全部答案。
\end{zhengju}

请留意这个故事的结局有多彻底：在漫长的共同生活里，原本属于细菌的大部分基因，要么丢失了，要么转移到了宿主的细胞核里，相应的蛋白质再由细胞质制造、运回线粒体。自由细菌退化成了“器官”，宿主细胞再也离不开它。竞争的舞台从“两个物种之间”转移到了“一个更大的个体内部”——原来的两个伙伴，合并成了一个更高层次的新个体。这件事给了我们一个重要的提示：\textbf{生命史上最大的几次创新，不是某个零件变厉害了，而是原本独立的复制单元合并成了新的整体}。

\section{选择到底在哪一层发生}

现在可以正面回答本章的核心问题了。把生命的复制单元按“谁套着谁”排一排：

\begin{center}
\begin{tikzpicture}[scale=1]
\draw[thick] (0,0) ellipse (4.6 and 3.0);
\draw[thick] (0,0) ellipse (3.2 and 2.0);
\draw[thick] (0,-0.1) ellipse (1.9 and 1.1);
\draw[thick] (0,-0.2) ellipse (0.8 and 0.45);
\node at (0,2.55) {群体};
\node at (0,1.55) {个体};
\node at (0,0.55) {细胞};
\node at (0,-0.2) {基因};
\end{tikzpicture}
\end{center}

（这是人为的表示法：真实的群体、个体、细胞当然不是这样整齐的同心圆，嵌套关系也常比图画得更复杂。）

\textbf{基因这一层}：选择偏爱“能让自己拷贝变多”的序列。上一章见过的转座元件就是纯粹在这一层成功的家伙——它们不为宿主做任何事，只会在基因组里复制粘贴自己，我们基因组里将近一半序列都是这类“搭便车者”留下的遗迹。只用“对个体有什么好处”解释不了它们，必须用基因层面的选择。

\textbf{个体这一层}：这是达尔文最熟悉的层次，也是前几章的主角。长脖子、保护色、强壮的腿——这些性状让携带它们的个体活得更久、繁殖得更多。

\textbf{群体这一层}：某些性状对个体不利、却对群体有利（比如节制繁殖、避免吃光资源），理论上，如果群体之间差异大、个体之间“拆伙重组”频繁，这样的性状也可能被保留。这叫\emph{群体选择}。要坦白告诉你：群体选择在自然界里到底有多重要，学界争论了几十年，至今没有完全一致的结论。多数研究者认为它在狭窄的条件下能起作用，但不能随手拿来解释一切——这是还没有标准答案的前沿，不是教科书遗漏了答案。

关键点在于：这三层选择\textbf{同时发生，方向可以相反}。一个基因层面的赢家（转座子）可能是个体层面的负担；一个个体层面的输家（蜇人的工蜂）可能是基因层面的大赢家。所以“这个性状为什么存在”从来不是一个问题，而是至少三个问题——必须先问清楚“在哪一层算账”。这就是本章标题里“多层次选择”的意思。

\begin{wuqu}
“《自私的基因》这本书说明：生物天生自私，而且每种行为的背后都有某一个‘自私基因’在操控。”——这是双重的误读。第一，“自私的基因”只是比喻：基因没有动机，不会“想”做任何事；这个说法的真实含义是，演化账本上最稳定的记账单位是能代代复制的序列。第二，恰恰是这套“基因自私”的逻辑，预言了合作——亲缘选择、互惠、共生，全都是“自私的复制单元”在特定条件下走出来的出路，蜜蜂和吸血蝙蝠就是证据。第三，真实的行为几乎从不是“一个基因一个行为”：行为由成百上千个基因、发育过程、经验和环境共同塑造，把每种行为都编成“某个基因的策略故事”，是不负责任的简化。还有一种更隐蔽的同类错误：见到任何性状都认定它“一定有适应意义”。第 77 章讲过漂变和中性演化——有些特征只是偶然的产物，或者是别的适应特征的副产品。负责任的演化解释必须给出能被检验的预测，否则就只是讲故事。
\end{wuqu}

\section{作弊者与警察：多细胞身体的内部治理}

理解了多层次选择，就能看懂生命史上最惊险的一跃：\textbf{从单细胞到多细胞}。

这件事在演化中独立发生过许多次，而每一次都要解决同一个难题：身体里的每个细胞都保留着分裂的能力，凭什么不“单干”？一个只顾自己疯狂分裂、把能量都用来复制自己的细胞，短期内是体内竞争的赢家——可如果所有细胞都这么干，多细胞个体就散架了。用这一章的语言说：\textbf{细胞层面的选择偏爱作弊，个体层面的选择惩罚作弊}，多细胞生命要存在，就必须让上层压制住下层。

你的身体为此配备了一整套“内部治理”机制：

\begin{itemize}[itemsep=2pt]
\item \textbf{单细胞瓶颈}：每一代都从一个受精卵重新开始。上一章讲过胚胎发育——整个身体的几十万亿个细胞，全都来自同一个细胞的克隆扩增，基因几乎完全相同。这意味着体内细胞之间几乎是“一家人”，亲缘系数接近 1，细胞层面的竞争失去了演化动力：帮身体就是帮自己基因的唯一出路。而且任何在体细胞里出现的新突变，都进不了下一代。
\item \textbf{生殖系与体细胞分工}：只有生殖细胞（将来变成精子、卵子的那一系）能把基因传下去，其余体细胞无论多“成功”都注定绝后。这就像蜂巢里的工蜂：体细胞是为生殖系打工的“不育品级”。
\item \textbf{执法系统}：细胞内部有监视 DNA 损伤的蛋白质（比如著名的 p53），发现修复不了的损伤就启动凋亡程序，让细胞有序自杀；组织里有“接触抑制”，细胞挤到一定程度就停止分裂；免疫系统则不断巡逻，清除行为异常的细胞。
\end{itemize}

把这三条放在一起看，你会意识到一件惊人的事：\textbf{你的身体本身就是一个演化出来的“社会”}，不育、自杀、警察，这些机制全是用来把“细胞层面自私的赢家”扼杀在摇篮里。而这套防御体系偶尔失效的时候，就是癌症。

\section{癌症：发生在身体里的演化}

癌症通常被描述为“细胞的叛变”，但用这一章的工具，可以说得更准确：\textbf{癌症是体细胞群体在身体内的一场演化}。把演化的三要素逐个对号入座：

\begin{itemize}[itemsep=2pt]
\item \textbf{变异}：每次细胞分裂复制 DNA 时都可能出错，加上紫外线、烟草里的化学物质（第 33 章讲过的致癌物怎样损伤 DNA）等因素，体细胞不断积累突变。
\item \textbf{选择}：绝大多数突变无害或有害，但极少数突变恰好踩中关键开关——让细胞无视“停止分裂”的信号、逃避凋亡、骗过免疫监视。携带这些突变的细胞分裂得比邻居快，它们的克隆就扩张。肿瘤内部不是铁板一块，而是许多突变各不相同的细胞克隆在\textbf{竞争}养分和空间。
\item \textbf{遗传}：每个细胞把突变传给子细胞，优势克隆的特征被继承和放大。
\end{itemize}

这也解释了一个临床上反复出现的难题：化疗或靶向药杀死了肿瘤里 99\% 的细胞，幸存的少数细胞却恰好携带耐药性突变，它们的克隆卷土重来——这和细菌对抗生素产生抗性（第 77 章）是\emph{同一套}演化逻辑，只是舞台换到了身体里。正因为癌症是演化过程，近年来一个方向是把演化思维引入治疗研究：例如不追求一次用最大剂量杀光肿瘤细胞，而是调整用药节奏，让敏感细胞“看住”耐药细胞、延缓耐药克隆的扩张（这类思路叫适应性治疗）。要强调：这些都还在临床试验阶段，证据等级各不相同，具体到任何个人的治疗选择，只能由医生和患者依据当时的证据决定——科普读物到此为止。

先来感受一下“为什么防御系统必须这么严密”——用数量级说话：

\begin{qiaoliang}
你的基因组约有 $3\times 10^{9}$ 个碱基对。体细胞每分裂一次，DNA 复制酶要抄完这 30 亿个字母，错误率大约是每复制一个碱基 $10^{-10}$ 到 $10^{-9}$ 次出一次错（复制后还有校对修复，第 66 章讲过这套纠错机制）。按这个数量级，\textbf{每次细胞分裂平均引入约 1 个新的点突变}。你的骨髓、肠道上皮、皮肤每天都在大量分裂，全身一天的细胞分裂次数约为 $10^{12}$ 量级。也就是说，仅仅今天一天，你身体里就产生了大约 $10^{12}$ 个携带新突变的细胞——一年下来，几乎所有可能的单点突变都在你身体的某个角落里出现过。这样算完，问题就不是“为什么有人会得癌症”，而是“\textbf{为什么我们不是人人都早早得癌症}”。答案就是上一节那套层层设防：癌变需要同一个细胞接连积累多个特定突变，还要逃过凋亡和免疫监视，每一道防线都把概率压低几个数量级。癌症是防御体系罕见的失守，而不是身体的常态。
\end{qiaoliang}

\section{再看那罐蜂蜜}

回到开场的悬念。工蜂为什么放弃生育？因为蜂巢是一个典型的“超个体”：蜂王是生殖系，工蜂是体细胞。单细胞瓶颈（每群蜂始于一对蜂王与雄蜂的交配）、亲缘选择（巢里全是高度近亲的姐妹）、个体之上的新层次（整群蜂一起过冬、一起御敌）——前面用来解释你身体的每一条逻辑，在蜂巢里都有一一对应的位置。工蜂的螫针为什么敢带着倒钩、蜇人即死？因为决定这根针命运的“账本”不在蜇人的那只工蜂身上：它死后，巢里成千上万携带相同基因的姐妹活了下来。倒钩甚至是件“好设计”——留在皮肤里的针连同毒囊会继续搏动、注入更多毒液，对蜂巢的保护更彻底。在个体层面这是自杀，在基因层面这是高分答卷。

而那罐蜂蜜本身，是这个“超个体”的集体储蓄：没有一只蜜蜂尝过自己劳动的全部果实，也没有任何一只蜜蜂“打算”为群体奉献——蜜蜂没有这种念头。整套秩序是从“复制单元的差异命运”这条冷冰冰的统计规则里长出来的。这跟第 26 章讲过的道理一脉相承：分子不“想要”稳定，只是低能量状态在统计上占优势；蜜蜂不“想要”牺牲，只是携带这些行为基因的蜂群在统计上留下了更多后代。拟人化是入门的拐杖，能量与统计才是地板。

\section{本篇的终点，下一篇的起点}

这一篇从化石和家谱出发，走到变异、选择、漂变、物种形成、基因组的自我改造，最后停在“选择发生在哪一层”这个问题上。你已经看到，演化的故事从不只是“强者生存”四个字：它有合作与背叛、有层次与博弈、有争议与未知，而且每一层都可以追溯到更底层的化学——根瘤菌耗掉的 ATP、固氮酶拆开的氮氮三键、复制酶抄错的那个碱基。

但请注意这一章里反复出现的一类句子：“把 DNA 序列拿来比对，发现线粒体和某类细菌最接近。”——我们凭什么能“读”出几十亿个碱基？读出来以后，能不能改？改了以后，能不能治病、能不能育种、会不会带来新的风险？下一篇《读写生命与共同未来》就从这里出发：第 81 章讲怎样读出 DNA 序列，第 82 章讲怎样改写基因，然后一路走向基因治疗、农业与生态的抉择。读懂这本书的人，迟早要参与书写它。

\section*{练一练}

\begin{sikaoti}
\item 画一张信息流图：以蜂巢为对象，标出蜂王、工蜂、雄蜂之间“基因从谁流向谁”。用图说明：为什么工蜂不育，却仍然能让自己携带的基因进入下一代？
\item 一只动物牺牲自己可以救亲属。若救的是亲兄弟姐妹（$r=1/2$），按汉密尔顿法则，至少救活几个才可能在演化上“够本”？若救的是侄子侄女（$r=1/4$）呢？这个计算做了哪些简化假设？
\item 画一张根瘤里的物质流图：标出糖和含氮物质分别在植物与根瘤菌之间朝哪个方向流动，并在旁边注明驱动这两条流各自的能量来源（提示：一条流要靠 ATP）。
\item 把一个正在生长的肿瘤当作一个“演化中的群体”，画出它的变异—选择—遗传循环图，并在图上标出：一次化疗杀伤 99\% 细胞之后，剩下的细胞会发生什么？这和农田里反复使用同一种除草剂有什么相似之处？
\item “孔雀的长尾巴是为了让物种更好看”“眼泪是为了排出压力激素”——这类解释犯了本章批评的什么错误？任选一句，写出一个能让它变成合格科学假说的、可被检验的预测。
\item 用本章的嵌套层次图解释：为什么“鹿跑得快是为了整个物种不被狼吃光”这种说法，比“跑得快的那只鹿活了下来”要可疑得多？在什么额外条件下，“为群体好”的解释才可能成立？
\end{sikaoti}
